\section*{NeurIPS Paper Checklist}

\begin{enumerate}

\item {\bf Claims}
    \item[] Question: Do the main claims made in the abstract and introduction accurately reflect the paper's contributions and scope?
    \item[] Answer: \answerYes{}
    \item[] Justification: The abstract and introduction center the graded-visibility study and its corrected-RGB follow-up. They distinguish the development cohort, the 25 reused central identities, the 75 newly captured extension identities, and the prospective 48-training/48-test split, history-matched from capacity-matched comparison, and linear and small nonlinear readouts from causal memory or absent model state. The experiments report central results, fresh-cohort refitting, and fixed-probe transfer separately.


\item {\bf Limitations}
    \item[] Question: Does the paper discuss the limitations of the work performed by the authors?
    \item[] Answer: \answerYes{}
    \item[] Justification: The interpretation subsection discusses the limited settings, reuse of central identities in the RGB follow-up, material appearance changes, condition-specific probes, finite reset windows, unequal estimator capacity, and conditional bootstrap uncertainty. Historical BGR observations are kept separate from the corrected-RGB results.


\item {\bf Theory assumptions and proofs}
    \item[] Question: For each theoretical result, does the paper provide the full set of assumptions and a complete (and correct) proof?
    \item[] Answer: \answerNA{}
    \item[] Justification: The paper presents an empirical study and does not introduce theoretical results, theorems, or formal proofs.


    \item {\bf Experimental result reproducibility}
    \item[] Question: Does the paper fully disclose all the information needed to reproduce the main experimental results of the paper to the extent that it affects the main claims and/or conclusions of the paper (regardless of whether the code and data are provided or not)?
    \item[] Answer: \answerYes{}
    \item[] Justification: The methods and appendix disclose preprocessing, checkpoint identity, alignment kernels, all fold memberships in the accompanying configuration, alpha selection, model inputs, and bootstrap settings. The recovered historical predictions and exact targets were verified. The accompanying compact snapshot includes new held-out predictions and a tested CPU command that reconstructs the statistics and figures without inference.


\item {\bf Open access to data and code}
    \item[] Question: Does the paper provide open access to the data and code, with sufficient instructions to faithfully reproduce the main experimental results, as described in supplemental material?
    \item[] Answer: \answerNo{}
    \item[] Justification: A complete local statistical-reproduction snapshot accompanies this revision, including predictions, targets, folds, source, provenance, and matched example images the fresh 75-identity extension, and the prospective variable-speed diagnostic. The complete original capture corpus was recovered and verified. Publication of the new snapshot and verification of anonymous public access are still pending; a project link alone does not establish that the new artifacts are publicly available.


\item {\bf Experimental setting/details}
    \item[] Question: Does the paper specify all the training and test details (e.g., data splits, hyperparameters, how they were chosen, type of optimizer) necessary to understand the results?
    \item[] Answer: \answerYes{}
    \item[] Justification: Sections 3.1--3.6 describe cohort construction, the fixed primary representation, grouped splits, ridge alpha selection, preprocessing, baseline inputs, and inferential units. Section 4 separates the central study from the earlier development protocol.


\item {\bf Experiment statistical significance}
    \item[] Question: Does the paper report error bars suitably and correctly defined or other appropriate information about the statistical significance of the experiments?
    \item[] Answer: \answerYes{}
    \item[] Justification: Section 3.5 defines paired 95 percent percentile intervals from 10,000 whole-identity resamples and explains that they condition on fitted probes and are not simultaneous bounds. Figures and tables distinguish mean identity RMSE from single-age absolute error. New comparisons are estimation analyses; historical Holm-corrected tests are not relabeled as confirmatory evidence for the follow-up.


\item {\bf Experiments compute resources}
    \item[] Question: For each experiment, does the paper provide sufficient information on the computer resources (type of compute workers, memory, time of execution) needed to reproduce the experiments?
    \item[] Answer: \answerNo{}
    \item[] Justification: The new follow-up reports GPU, CPU, memory, software, measured extraction and fitting times, and preparation time separately. Recovered receipts identify the historical NVIDIA L4 extraction environment. Complete total-project and earlier-development costs are not established, so this answer remains No for the full history rather than implying all historical compute was measured.


\item {\bf Code of ethics}
    \item[] Question: Does the research conducted in the paper conform, in every respect, with the NeurIPS Code of Ethics \url{https://neurips.cc/public/EthicsGuidelines}?
    \item[] Answer: \answerYes{}
    \item[] Justification: The study uses a controlled simulated environment and does not involve human subjects, personal data, or scraped content. The research and anonymous submission follow the NeurIPS Code of Ethics.


\item {\bf Broader impacts}
    \item[] Question: Does the paper discuss both potential positive societal impacts and negative societal impacts of the work performed?
    \item[] Answer: \answerNA{}
    \item[] Justification: This is a diagnostic study in a simulated environment, not a deployed decision system. Its immediate benefit is clearer evaluation of world-model readouts; extrapolation to reliable physical decision-making is explicitly outside the measured claims.


\item {\bf Safeguards}
    \item[] Question: Does the paper describe safeguards that have been put in place for responsible release of data or models that have a high risk for misuse (e.g., pre-trained language models, image generators, or scraped datasets)?
    \item[] Answer: \answerNA{}
    \item[] Justification: The released benchmark and evaluation materials do not introduce a new high-risk generative model or contain personal or human-subject data.


\item {\bf Licenses for existing assets}
    \item[] Question: Are the creators or original owners of assets (e.g., code, data, models), used in the paper, properly credited and are the license and terms of use explicitly mentioned and properly respected?
    \item[] Answer: \answerYes{}
    \item[] Justification: MineWorld \texttt{700M\_16f} is credited and used under the MIT License, and Fabric API \texttt{0.42.0+1.16} is used under Apache-2.0. Minecraft Java Edition 1.16.5 is used without redistributing game files and in accordance with the Mojang/Microsoft EULA and Usage Guidelines.


\newpage
\item {\bf New assets}
    \item[] Question: Are new assets introduced in the paper well documented and is the documentation provided alongside the assets?
    \item[] Answer: \answerYes{}
    \item[] Justification: The accompanying snapshot documents file schemas, model/input configurations, bootstrap settings, checksums, exact folds, execution commands, and the distinction between statistical reconstruction and full GPU reproduction. Public anonymous distribution remains a separate pending release step, as stated above.


\item {\bf Crowdsourcing and research with human subjects}
    \item[] Question: For crowdsourcing experiments and research with human subjects, does the paper include the full text of instructions given to participants and screenshots, if applicable, as well as details about compensation (if any)? 
    \item[] Answer: \answerNA{}
    \item[] Justification: The study does not involve crowdsourcing, human participants, or data collected from human subjects.


\item {\bf Institutional review board (IRB) approvals or equivalent for research with human subjects}
    \item[] Question: Does the paper describe potential risks incurred by study participants, whether such risks were disclosed to the subjects, and whether Institutional Review Board (IRB) approvals (or an equivalent approval/review based on the requirements of your country or institution) were obtained?
    \item[] Answer: \answerNA{}
    \item[] Justification: The study does not involve human participants or human-subject data, so IRB approval or an equivalent review was not required.


\item {\bf Declaration of LLM usage}
    \item[] Question: Does the paper describe the usage of LLMs if it is an important, original, or non-standard component of the core methods in this research? Note that if the LLM is used only for writing, editing, or formatting purposes and does \emph{not} impact the core methodology, scientific rigor, or originality of the research, declaration is not required.
    \item[] Answer: \answerYes{}
    \item[] Justification: LLM assistance supported the prospective diagnostic design, implementation, analysis scripts, debugging, and manuscript preparation, as disclosed in the diagnostic appendix. It did not generate the capture images, targets, or reported numerical findings. Measurements are bound to executed captures and model predictions with explicit validation checks.

\end{enumerate}
